\documentclass[11pt]{article}
\usepackage[a4paper,margin=21mm]{geometry}
\usepackage{fontspec}
\setmainfont{Latin Modern Roman}
\newfontfamily\CJKfont{Noto Serif CJK SC}
\XeTeXlinebreaklocale "zh"
\XeTeXlinebreakskip=0pt plus 1pt
\usepackage{amsmath,amssymb,amsthm,amscd,graphicx,color}
\usepackage{xurl}
\usepackage[unicode,hidelinks]{hyperref}
\allowdisplaybreaks
\setlength\emergencystretch{3em}
\numberwithin{equation}{section}

\numberwithin{equation}{section}

\newtheorem{Theorem}{Theorem}[section]
\newtheorem{Corollary}[Theorem]{Corollary}
\newtheorem{Lemma}[Theorem]{Lemma}
\newtheorem{Proposition}[Theorem]{Proposition}
\newtheorem{prop}[Theorem]{Result}
 { \theoremstyle{definition}
\newtheorem{Definition}[Theorem]{Definition}
\newtheorem{Note}[Theorem]{Note}
\newtheorem{Example}[Theorem]{Example}
\newtheorem{Remark}[Theorem]{Remark} }

\newcommand{\FA}{\mathfrak{A}}
\newcommand{\FB}{\mathfrak{B}}
\newcommand{\FC}{\mathfrak{C}}
\newcommand{\FD}{\mathfrak{D}}
\newcommand{\FE}{\mathfrak{E}}
\newcommand{\FH}{\mathfrak{H}}
\newcommand{\FK}{\mathfrak{K}}
\newcommand{\FL}{\mathfrak{L}}
\newcommand{\FR}{\mathfrak{R}}
\newcommand{\Fa}{\mathfrak{a}}
\newcommand{\Fd}{\mathfrak{d}}
\newcommand{\Fg}{\mathfrak{g}}
\newcommand{\Fh}{\mathfrak{h}}
\newcommand{\Fk}{\mathfrak{k}}
\newcommand{\Fl}{\mathfrak{l}}
\newcommand{\FG}{\mathfrak{G}}
\newcommand{\FX}{\mathfrak{X}}
\newcommand{\FY}{\mathfrak{Y}}
\newcommand{\Fn}{\mathfrak{n}}
\newcommand{\Ff}{\mathfrak{f}}
\newcommand{\Frakc}{\mathfrak{c}}
\newcommand{\Fp}{\mathfrak{p}}
\newcommand{\Fq}{\mathfrak{q}}
\newcommand{\Fs}{\mathfrak{s}}
\newcommand{\FZ}{\mathfrak{Z}}
\newcommand{\FN}{\mathfrak{N}}
\newcommand{\Kc}{\mathcal{K}}
\newcommand{\Hc}{\mathcal{H}}
\newcommand{\Cc}{\mathcal{C}}
\newcommand{\Tc}{\mathcal{T}}
\newcommand{\Mc}{\mathcal{M}}
\newcommand{\Zc}{\mathcal{Z}}
\newcommand{\Fc}{\mathcal{F}}
\newcommand{\Oc}{\mathcal{O}}
\newcommand{\Ac}{\mathcal{A}}
\newcommand{\Xc}{\mathcal{X}}

\newcommand{\hmod}{{_H\mathcal{M}}}
\newcommand{\amod}{{_A\mathcal{M}}}
\newcommand{\Bmod}{{_B\mathcal{M}}}
\newcommand{\cmod}{{_C\mathcal{M}}}
\newcommand{\dmod}{{_D\mathcal{M}}}

\newcommand{\kmod}{{_K\mathcal{M}}}
\newcommand{\hhmod}{{_{H^2}\mathcal{M}}}
\newcommand{\md}{\mathbb{M}\text{d}}
\newcommand{\Dc}{\mathcal{D}}
\newcommand{\Nc}{\mathcal{N}}
\newcommand{\Gc}{\mathcal{G}}
\newcommand{\Ec}{\mathcal{E}}
\newcommand{\aydc}{\widehat{a\mathcal{YD}}}
\newcommand{\aydm}{a\mathcal{YD}}
\newcommand{\indlimD}{\displaystyle{\varinjlim_{\Delta}}}
\newcommand{\projlimD}{\displaystyle{\varprojlim_{\Delta}}}
\newcommand{\indlimlinf}{\displaystyle{\varinjlim_{\Lambda_\infty}}}
\newcommand{\projlimlinf}{\displaystyle{\varprojlim_{\Lambda_\infty}}}
\newcommand{\indliml}{\displaystyle{\varinjlim_{\Lambda}}}
\newcommand{\projliml}{\displaystyle{\varprojlim_{\Lambda}}}

\newcommand{\ns}[1]{~\hspace{-3pt}_{\left<#1\right>}}
\newcommand{\ps}[1]{~\hspace{-3pt}^{^{\left(#1\right)}}}
\newcommand{\sns}[1]{~\hspace{-3pt}_{\left<\overline{#1}\right>}}
\newcommand{\nsb}[1]{~\hspace{-3pt}_{^{[#1]}}}
\newcommand{\bb}[1]{{}_{\underline{#1}}}
\newcommand{\wid}[1]{\dot{\overbrace{#1}}}
\newcommand{\lu}[1]{\;^{#1}\hspace{-2pt}}
\newcommand{\li}[1]{\;_{#1}\hspace{-1pt}}
\newcommand{\vect}{\text{Vec}}
\newcommand{\id}{1}
\newcommand{\bimod}{\mathbb{B}\text{md}}
\newcommand{\bmc}{\mathcal{BMC}}
\newcommand{\Hm}{\mathcal{H}om}
\newcommand{\Md}{\mathcal{M}od}
\newcommand{\la}{\triangleright}
\newcommand{\ra}{\triangleleft}
\newcommand{\lba}{\blacktriangleright}
\newcommand{\rba}{\blacktriangleleft}
\newcommand{\Bc}{\mathcal{B}}
\newcommand{\dF}{{^*F}}
\newcommand{\bt}{\boxtimes}
\newcommand{\mmm}{{\text{-mod}}}

\newcommand{\bx}{\boxtimes}

\newcommand{\ot}{\otimes}

\newcommand{\zp}{\mathbb{Z}/p}
\usepackage{mathtools}\mathtoolsset{showonlyrefs=true}
\makeatletter\newlabel{appx}{{A}{}{Original source locator}{}{}}\makeatother
\begin{document}
\begin{center}\Large Mixed anti-Yetter-Drinfeld contramodules and mixed complexes of stable contramodules for Sweedler Hopf algebra are not DG-equivalent over algebraically closed characteristic-zero fields\end{center}
\noindent\textbf{Stable ID:} 1343\quad\textbf{Author:} Ilya Shapiro\par
\noindent\textbf{Work:} Mixed vs Stable Anti-Yetter–Drinfeld Contramodules\par
\noindent\textbf{Source:} \url{https://sigma-journal.com/2021/026/}\par
\noindent\textbf{Version:} Official SIGMA 17 (2021) article 026, DOI 10.3842/SIGMA.2021.026; journal PDF and native TeX acquired 2026-09-30, exact bytes pinned by SHA256\par
\noindent\textbf{Rights:} CC BY 4.0. Authors retain copyright. \url{https://creativecommons.org/licenses/by/4.0/}. Reuse and typesetting adaptation; original language and mathematical content preserved.\par
\noindent\textbf{Difficulty (prior selection preserved):} {\CJKfont H2: The proof computes the twisted Hopf double and its central circle-action element, then decomposes the module problem into two central blocks. One block is identified with ordinary mixed complexes, while the nonsemisimple block gives two different differential graded algebras with distinct degree-minus-one Hochschild cohomology. The explicit nilpotent-block calculation supplies the obstruction to equivalence; the general invariance theorem alone does not decide it.}\par
\medskip\noindent\textbf{Editorial boundary note (not author text):} Selected primary Proposition3.6 for the four-dimensional Sweedler Hopf algebra over algebraically closed characteristic0 fields. Full mixed/stable DG-module definitions and twisted double/central sigma, precise DGAs, complete diagonalizable-block lemma, Taft multiplication/coproduct/antipode and dual-isomorphism formulas, grading/relations/sigma calculation and both p=2 central blocks are retained. The complete negative-degree Hochschild-cohomology formula, actual center/kernel computations and differing dimensions supply the new obstruction. Prior standard Taft self-duality is explicitly cited, with both isomorphism and inverse recipes and the literal reader-verification note retained; no generated verification. Semisimple/S-square identity and p>2/general-double comparisons are unselected; no supporting component receives an extra count. All author signs/grading/index formulas remain literal. Original SIGMA-class showonlyrefs equation-numbering behavior is reproduced with installed mathtools; no mathematical display text is altered.\par
\noindent\textbf{External original reference appx:} Original AppendixA/native448–487 compares familiar Drinfeld doubles; selected mixed-vs-stable proof works directly with the fully supplied twisted double2.1 and relations3.2 and does not require this unselected comparison.\par
\medskip\hrule\medskip\noindent\textbf{Author text begins below.}\par
\setcounter{section}{1}\setcounter{Theorem}{1}
% Original source lines 148--150
\begin{prop}
Let $H=T_2(-1)$, then the mixed complexes in the category of stable anti-Yetter--Drinfeld contramodules are not DG-equivalent to the category of mixed anti-Yetter--Drinfeld contramodules.
\end{prop}

\setcounter{section}{1}
% Original source lines 152--241
\noindent \textbf{Conventions.} All algebras $A$ in monoidal categories are assumed to be unital associative. Our~$H$ is a Hopf algebra over some fixed algebraically closed field $k$, of characteristic $0$, and $\vect$ denotes the category of $k$-vector spaces. For the purposes of this paper we are only interested in finite dimensional Hopf algebras. We~use the following version of Sweedler's notation: For~$h\in H$ we denote the coproduct $\Delta(h)\in H\ot H$ by $h^1\ot h^2$. The letter $S$ denotes the antipode of~$H$. The number $p$ is prime. Finally, DG stands for differential graded.

\section{Twisted Drinfeld double}
Let $H$ be a Hopf algebra. From~\cite{chern} we see that the study of the Hochschild and cyclic homologies of $\hmod$, the monoidal category of $H$-modules, reduces to the study of modules over a certain monad on $\hmod$. Recall that the consideration of Hochschild and cyclic homologies of monoidal categories is motivated by their recently discovered role~\cite{chern} in the understanding of Hopf-cyclic theory coefficients.

Briefly, we have the monad (see~\cite{chern} for more details):
\begin{gather}\label{themonad}
\mathop{\rm Hom}\nolimits_k(H,-)\colon \ \hmod\to\hmod
\end{gather}
with the $H$-module structure on $\mathop{\rm Hom}\nolimits_k(H,V)$:
\begin{gather}
x\cdot\varphi=x^2\varphi\big(S\big(x^3\big)(-)x^1\big),
\end{gather}
for $x\in H$ and $\varphi\in \mathop{\rm Hom}\nolimits_k(H,V)$. The unit $1_V\colon {\rm Id}(V)\to \mathop{\rm Hom}\nolimits_k(H,V)$ is
\begin{gather}
1_V(v)(h)=\epsilon(h)v
\end{gather}
and a crucial central element $\sigma_V\colon {\rm Id}(V)\to \mathop{\rm Hom}\nolimits_k(H,V)$ is
\begin{gather}
\sigma_V(v)(h)=hv.
\end{gather}

The anti-Yetter--Drinfeld contramodules then coincide with modules over this monad (shown in~\cite{chern}), while the stable ones consist of those for which the action of $\sigma$ agrees with that of $1$, and the mixed ones introduced in~\cite{chern} are the homotopic version of this on the nose requirement.

Recall the mixed complexes of~\cite{kassel}. These are complexes of vector spaces $(V^\bullet,d)$ with a~homo\-topy $h$ such that $dh+hd=0$. We~can replace vector spaces with $R$-modules for some ring~$R$. Note that as observed in~\cite{kassel}, the DG-category of mixed complexes in $R$-modules is isomorphic to the DG-category of DG-modules over $R[\theta]$, where $\theta$ is a freely and centrally adjoined degree~$-1$~graded commutative variable (naturally $R$ itself is placed in degree $0$, so that $R[\theta]=R\to R$ as a~comp\-lex) and $d=0$ on $R[\theta]$. The action of $\theta$ gives the homotopy $h$.

We can generalize the considerations of~\cite{kassel} so as to apply to our particular situation. Namely, let $z\in Z(R)$, i.e., $zr=rz$ for all $r\in R$. Define a DG-algebra $R[\theta]$ by placing $R$ in degree $0$ and $\theta$ in degree $-1$. Let $\theta$ commute with $R$ and itself, so in particular $\theta^2=0$. So far it is as above. Now define the differential to be $0$ on $R$ and $d\theta=z$. This is well defined and unique by the Leibniz rule. We~observe that the category of DG-modules over $R[\theta]$ consists of complexes of $R$-modules equipped with a homotopy $h$ such that $dh+hd=z$.

Now recall from~\cite{chern}:

\begin{Definition}
We say that $(M^\bullet,d,h)$ is a mixed anti-Yetter--Drinfeld contramodule if $(M^\bullet,d)$ is a complex of contramodules, i.e., modules over the monad \eqref{themonad}, and $h$ is a homotopy annihilating $\sigma-1$. More precisely,
\begin{gather}
dh+hd=\sigma-1.
\end{gather}
\end{Definition}

In this section we will define an explicit DG-algebra that will yield the mixed anti-Yetter--Drinfeld (aYD) contramodules (for $H$ finite dimensional) as its DG-modules. The construction of the twisted convolution algebra below is analogous to the classical Drinfeld double $D(H)$ and its anti-version $D_a(H)$~\cite{HKRS2} (we review these in Appendix~\ref{appx}, where we expand upon this comparison).

\begin{Definition}\label{dhhat}
Let $H$ be a Hopf algebra with an invertible antipode $S$, define a twisted double~$\widehat{D}(H)$ as follows. The multiplication on $\widehat{D}(H):={\rm End}(H)$ is
\begin{gather}\label{mult}
(f\star g)(h)= f\big(h^1\big)^2g\big(S\big(f\big(h^1\big)^3\big)h^2f\big(h^1\big)^1\big),
\end{gather}
thus the multiplicative identity, which we denote by $1$ is $\epsilon(-)1$, and the central element $\sigma(h)=h$ is invertible with inverse $S^{-1}$.
\end{Definition}

Definition~\ref{dhhat} is extracted from the monad~\eqref{themonad} with the sole purpose consisting of making the following lemma a tautology.

\begin{Lemma}\label{bas}
Let $H$ be a finite dimensional Hopf algebra.
\begin{itemize}\itemsep=0pt
\item The category of anti-Yetter--Drinfeld contramodules over $H$ is isomorphic to $\widehat{D}(H)$-mo\-dules.

\item The category of stable anti-Yetter--Drinfeld contramodules is isomorphic to modules over
\begin{gather}
A:=\widehat{D}(H)/(\sigma-1).
\end{gather}

\item The DG-category of mixed anti-Yetter--Drinfeld contramodules is isomorphic to DG-mo\-dules over the DG algebra
\begin{gather}
B:=\widehat{D}(H)[\theta],
\end{gather}
where $\theta$ is a freely adjoined degree $-1$ graded commutative variable and $d\theta=\sigma-1$, with $d|_{\widehat{D}(H)}=0$.
\end{itemize}
\end{Lemma}
\begin{proof}
In the finite dimensional case, as vector spaces, $\widehat{D}(H)={\rm End}(H)\simeq H^*\ot H$. Furthermore, as an algebra, $\widehat{D}(H)$ is the quotient of the free product algebra, generated by $H^*$ and $H$, by the relation:
\begin{gather}
\label{mult1}
h \chi= \chi\big(S\big(h^3\big)(-)h^1\big)h^2,
\end{gather}
where $h\in H$, $\chi\in H^*$. Thus, modules over the algebra are both $H$-modules and $H$-cont\-ra\-mo\-du\-les (same as $H^*$-modules for $H$ finite dimensional). The two actions satisfy the requisite compatibility condition for contramodules, as specified in~\cite{contra}, and ensured by \eqref{mult1}. Clearly, modules over $A$ consist of the full subcategory of objects on which $\sigma$ acts by identity, these are exactly the stable contramodules. Finally, a DG-module over $B$ is just a complex of $\widehat{D}(H)$-modules with a homotopy given by the action of $\theta$. The condition $d\theta=\sigma-1$ ensures that $dh+hd=\sigma-1$ on~$M^\bullet$.
\end{proof}

Our main goal is to compare the category of mixed
aYD contramodules to the category of~mixed complexes of stable aYD
contramodules. By the preceding lemma this means determining when,
and more interestingly when not, the category of DG-modules over $B$ is DG-equivalent to $A[\theta]$-modules (with $\theta$ of degree $-1$ and $d=0$). The study of Hopf-cyclic cohomology has thus far only concerned itself with the latter.

The following simple lemma takes care of a lot of cases.

\begin{Lemma}\label{diag}
Let $H$ be a finite dimensional Hopf algebra and suppose that the action of $\sigma-1$ on~$\widehat{D}(H)$ is diagonalizable. Then the categories of mixed complexes in stable $aYD$-contramodules and mixed $aYD$-contramodules are $DG$-equivalent.
\end{Lemma}

\begin{proof}
Since the action of the central element $\sigma-1$ on $\widehat{D}(H)$ is diagonalizable we may decom\-pose $\widehat{D}(H)$ as a product of algebras $\widehat{D}(H)_0\oplus \widehat{D}(H)_+$ with $\widehat{D}(H)_0$ being the $0$-eigenspace and~$\widehat{D}(H)_+$ all the other eigenspaces. We~have an inclusion of DGAs: $\widehat{D}(H)_0[\theta]\to B$ that induces an isomorphism on cohomology. Namely, as complexes $B$ is $\widehat{D}(H)\stackrel{\sigma-1}{\to}\widehat{D}(H)$ whereas $\widehat{D}(H)_0[\theta]$ is $\widehat{D}(H)_0\stackrel{0}{\to}\widehat{D}(H)_0$ and $\sigma-1$ is invertible on $\widehat{D}(H)_+$. Note that $\widehat{D}(H)_0\simeq A$ and we are done.
\end{proof}



% Original source lines 251--337
\section{Taft Hopf algebras}\label{sec:blah}
We fix a prime $p$ and a primitive $p$th root of unity $\xi\in k$ in the following. The Taft Hopf algebra $T_p(\xi)$~\cite{taft} is generated as a $k$-algebra by $g$ and $x$ with the relations \begin{gather}
g^p=1,\qquad
x^p=0,
\\
gx=\xi xg.
\end{gather}
It is sometimes called the quantum ${\mathfrak{sl}}_2$ Borel algebra. It is $p^2$ dimensional over $k$. Furthermore, the coalgebra structure is
\begin{gather}
\Delta(g)=g\ot g,\qquad \Delta(x)=x\ot 1+g\ot x
\end{gather}
with $\epsilon(g)=1$, $\epsilon(x)=0$, and thus $S(g)=g^{-1}$, while $S(x)=-g^{-1}x$. Note that
\begin{gather}
S^2(x)=\xi^{-1}x\neq x,
\end{gather}
making $T_2(-1)$ the smallest Hopf algebra with $S^2\neq {\rm Id}$. The Taft algebra $T_2(-1)$ is somewhat different from the other $T_p(\xi)$ and has its own name: Sweedler’s Hopf algebra.

\subsection{The identification with the dual}
The Taft algebra is isomorphic to its dual. We~need some explicit formulas establishing the isomorphism $T_p(\xi)\simeq T_p(\xi)^*$ and its inverse. Suppose that $\omega$ is a $p$th root of unity, let
\begin{gather}
(n)_\omega=1+\cdots +\omega^{n-1}
\end{gather}
and
\begin{gather}
(n)_\omega !=(n)_\omega\cdots(1)_\omega.
\end{gather}

The verification of the following is left to the reader; key details can be found in~\cite{selfdual}. The lemma itself can be obtained from~\cite{nen}.
\begin{Lemma}\label{nenlem}
 As Hopf algebras
 \begin{gather}
 T_p(\xi)^*\simeq T_p(\xi).
 \end{gather}
\end{Lemma}
\begin{proof}
Consider a basis of $T_p(\xi)$: $\big\{g^ix^j\big\}_{ i,j=0}^{p-1}$ so that $\big\{\big(g^ix^j\big)^*\big\}$ denotes the dual basis of $T_p(\xi)^*$. Then the isomorphism of Hopf algebras and its inverse are given by
\begin{gather}
g^ix^j\mapsto(j)_{\xi^{-1}}!\sum_l\xi^{i(j+l)}\big(g^lx^j\big)^*
\end{gather}
and
\begin{gather}
(g^ix^j)^*\mapsto\dfrac{1}{p(j)_{\xi^{-1}}!}\sum_l\xi^{-l(i+j)}g^lx^j.
\end{gather}
\end{proof}

\begin{Corollary}\label{gens}The twisted double $\widehat{D}(T_p(\xi))$ is a quotient of $k\left\langle x,x',g,g'\right\rangle.$

\begin{itemize}\itemsep=0pt
\item The relations are
\begin{gather*}
x^p=x'^p=g^p-1=g'^p-1=0,
\\
gg'=g'g,\qquad
gx=\xi xg,\qquad
g'x'=\xi x'g',\qquad
gx'=\xi^{-1} x'g, \qquad
g'x=\xi^{-1} xg',
\\
xx'-\xi^{-1}x'x=1-\xi^{-1}g'^{-1}g.
\end{gather*}

\item The actions of $g'$ and $g$ on a $\widehat{D}(T_p(\xi))$-module $V$, yields a $(\mathbb{Z}/p)^2$-grading on $V$ by their eigenspaces, i.e., $g'$, $g$ act on $V_{ij}$ by $\xi^i$, $\xi^j$, respectively. Thus $x$ and $x'$ have degrees $(-1,1)$ and $(1,-1)$, respectively.

\item The $S^1$-action of $\sigma$ on $V_{ij}$ is
\begin{gather}\label{sigmaaction}
\sum_{l=0}^{p-1}\dfrac{\xi^{(i-l)(j+l)}}{(l)_{\xi^{-1}}!}x'^lx^l.
\end{gather}
\end{itemize}
\end{Corollary}

\begin{proof}
We use the identification of vector spaces $\widehat{D}(T_p(\xi))=(T_p(\xi))^*\ot T_p(\xi)$ as in Lemma~\ref{bas} followed by $(T_p(\xi))^*\ot T_p(\xi) \simeq T_p(\xi)\ot T_p(\xi)$ from Lemma~\ref{nenlem}. We~let
\begin{gather}
x'=x\ot 1, \qquad
x=1\ot x, \qquad
g'=g\ot 1, \qquad
g=1\ot g
\end{gather}
in the latter. To derive the rest of the relations we apply \eqref{mult1}. The action of $\sigma=\sum_{ij}\big(g^ix^j\big)^*\ot g^ix^j$ is computed on the graded components directly.
\end{proof}

Observe that it follows from Corollary~\ref{gens} that $gg'\in\widehat{D}(T_p(\xi))$ is central. Since we see that $(gg')^p=1$ so its action on $\widehat{D}(T_p(\xi))$ is diagonalizable with eigenvalues $\xi^s$, $s\in\zp$. Thus as an~algebra
\begin{gather}\label{split}
\widehat{D}(T_p(\xi))=\bigoplus_s \widehat{D}(T_p(\xi))/(gg'-\xi^s)
\end{gather}
so that it suffices to understand $\widehat{D}(T_p(\xi))/(gg'-\xi^s)$. There are two cases: $p=2$ and $p>2$. We~will begin by briefly discussing the latter (though without addressing the $S^1$-action), and then concentrate our attention on the former (with examining the $S^1$-action) to achieve the goal set out in the abstract.


\setcounter{subsection}{2}\setcounter{Theorem}{3}\setcounter{equation}{4}
% Original source lines 370--444
\subsection[The case of p=2]{The case of $\boldsymbol{p=2}$}
We need to describe the algebra $\widehat{D}(T_2(-1))$ in greater detail, paying particular attention to the element $\sigma$.

By \eqref{split} the category of DG-modules over $\widehat{D}(T_2(-1))[\theta]$ is a product of categories, $\Cc_0\times \Cc_1$, corresponding to the cases $s=0$ and $s=1$. We~will deal with both separately. More precisely, for a $\widehat{D}(T_2(-1))[\theta]$-module $V$, we decompose
\begin{gather}
V=(V_{00}\oplus V_{11})\oplus(V_{01}\oplus V_{10}).
\end{gather}

Observe that by the Corollary~\ref{gens} we have that
\begin{gather}
xx'+x'x=1+(-1)^s,
\end{gather}
so that
\begin{gather}\label{dd}
(x'x)^2=(1+(-1)^s)x'x.
\end{gather}
We see from \eqref{dd} that the minimal polynomial of $x'x$ depends only on $s$; this is exclusive to $p=2$ and makes this case tractable. Note that by \eqref{sigmaaction}:
\begin{gather}\label{sig}
\sigma|_{V_{00}}=1-x'x,\qquad \sigma|_{V_{11}}=-1+x'x,\qquad\text{and}\qquad \sigma|_{V_{01}}=\sigma|_{V_{10}}=1+x'x.
\end{gather} Let
\begin{gather}
D_s=\widehat{D}(T_2(-1))/(gg'-(-1)^s).
\end{gather}

We begin with $s=0$: The category of $D_0$-modules consists of $\mathbb{Z}/2$-graded vector spaces ($V=V_{00}\oplus V_{11}$) equipped with degree changing operators $x$ and $x'$ subject to the relation $xx'+x'x=2$. The action of $\sigma-1$ on $V_{00}$ is $-x'x$ and on $V_{11}$ it is $x'x-2$.

\begin{Lemma}
The category $\Cc_0$ consists of the mixed complexes of~{\rm \cite{kassel}}.
\end{Lemma}

 \begin{proof}
Note that $D_0/(\sigma-1)$-modules are just vector spaces. Indeed, let $y=x'/2$. For improved clarity, denote by $x_i$ and $y_i$ the action of $x$ and $y$, respectively, that originates at $V_{ii}$. After modding out by $\sigma-1$ we have by \eqref{sig} that $y_1x_0=0\implies x_1y_0=1$ and $y_0x_1=1$. So
\begin{gather}
y_0:V_{00}\simeq V_{11}: x_1.
\end{gather}
Furthermore, $x^2=y^2=0$ so that both $x_0$ and $y_1$ are the $0$ maps. Thus
\begin{gather}
V=V_{00}\oplus V_{11}\mapsto V_{00}
\end{gather}
establishes an equivalence of categories between $D_0/(\sigma-1)$-modules and $\vect$.

The action of $\sigma-1$ on $D_0$ is diagonalizable by \eqref{dd} and so by the proof of Lemma~\ref{diag} the algebras $D_0[\theta]$ ($d\theta=\sigma-1$) and $D_0/(\sigma-1)[\theta]$ ($d\theta=0$) are quasi-isomorphic. Thus, $\Cc_0$, being DG-equivalent to $D_0/(\sigma-1)[\theta]$-modules, is by the above discussion, equivalent to $k[\theta]$-modules. These are just the mixed complexes of~\cite{kassel}.
\end{proof}

 \begin{Remark}
 Note that not only does $\Cc_0$ consist of the usual mixed complexes but it also does not provide any evidence of the need for the mixed $aYD$-contramodules (see the proof above).
 \end{Remark}

Moving on to $s=1$ we find that things change for the better. Recall that $D_1$ is generated by
\begin{gather}
x,\ x',\ g
\end{gather}
subject to
\begin{gather}
x^2=x'^2=g^2-1=0, \qquad
xx'=-x'x, \qquad
gx=-xg, \qquad
gx'=-x'g.
\end{gather}
Furthermore, $\sigma-1$ acts as $x'x$ by \eqref{sig}.

\begin{Proposition}\label{notsscase}
Let $H=T_2(-1)$, then the mixed complexes in the category of stable anti-Yetter--Drinfeld contramodules are not DG-equivalent to the category of mixed anti-Yetter--Drinfeld contramodules.
\end{Proposition}
\begin{proof}
By the preceding discussion, i.e., the decomposition of the category into a product, it~suf\-fices to show that the categories $D_1[\theta]\text{-mod}$ (where $d\theta=x'x$) and $D_1/(x'x)[\theta]\text{-mod}$ (where $d\theta =0$) are not DG-equivalent.

Recall that the Hochschild cohomology $HH^i(C^\bullet)$ of a DG algebra $C^\bullet$ is an invariant of its DG category of modules~\cite{dginvar}. We~will compute $HH^{-1}$. In our simple case of a DG algebra $C=C^{-1}\stackrel{d}{\to} C^0$ concentrated in two degrees, we have
\begin{gather}
HH^{-1}(C)={\rm ker} \big(C^{-1}\stackrel{\alpha}{\to} C^0\oplus {\rm Hom}\big(C^0, C^{-1}\big)\big),
\end{gather}
where $\alpha(x)=(dx, [x,-])$.

The key observation here is that the center of $D_1$ is spanned by $1$, $xx'$, $xx'g$, while that of~$D_1/(x'x)$ is spanned by $1$. Thus, in the first case we get that $HH^{-1}$ is spanned by $xx'$ and~$xx'g$. In the second case it is spanned by $1$.
\end{proof}
\begin{thebibliography}{99}
\bibitem[5]{contra}
Brzezi\'nski T., Hopf-cyclic homology with contramodule coefficients, in
 Quantum groups and noncommutative spaces, \textit{Aspects Math.}, Vol.~E41,
 \href{https://doi.org/10.1007/978-3-8348-9831-9_1}{Vieweg + Teubner}, Wiesbaden, 2011, 1--8, \href{https://arxiv.org/abs/0806.0389}{arXiv:0806.0389}.

\bibitem[9]{HKRS2}
Hajac P.M., Khalkhali M., Rangipour B., Sommerh\"auser Y., Stable
 anti-{Y}etter--{D}rinfeld modules, \href{https://doi.org/10.1016/j.crma.2003.11.037}{\textit{C.~R.~Math. Acad. Sci. Paris}}
 \textbf{338} (2004), 587--590, \href{https://arxiv.org/abs/math.QA/0405005}{arXiv:math.QA/0405005}.

\bibitem[12]{kassel}
Kassel C., Cyclic homology, comodules, and mixed complexes, \href{https://doi.org/10.1016/0021-8693(87)90086-X}{\textit{J.~Algebra}}
 \textbf{107} (1987), 195--216.

\bibitem[16]{nen}
Nenciu A., Quasitriangular pointed {H}opf algebras constructed by {O}re
 extensions, \href{https://doi.org/10.1023/B:ALGE.0000026785.03997.60}{\textit{Algebr. Represent. Theory}} \textbf{7} (2004), 159--172.

\bibitem[18]{selfdual}
Radford D.E., Westreich S., Trace-like functionals on the double of the {T}aft
 {H}opf algebra, \href{https://doi.org/10.1016/j.jalgebra.2004.04.023}{\textit{J.~Algebra}} \textbf{301} (2006), 1--34.

\bibitem[20]{chern}
Shapiro I., Categorified {C}hern character and cyclic cohomology,
 \href{https://arxiv.org/abs/1904.04230}{arXiv:1904.04230}.

\bibitem[21]{taft}
Taft E.J., The order of the antipode of finite-dimensional {H}opf algebra,
 \href{https://doi.org/10.1073/pnas.68.11.2631}{\textit{Proc. Nat. Acad. Sci. USA}} \textbf{68} (1971), 2631--2633.

\bibitem[22]{dginvar}
To\"en B., The homotopy theory of {$dg$}-categories and derived {M}orita
 theory, \href{https://doi.org/10.1007/s00222-006-0025-y}{\textit{Invent. Math.}} \textbf{167} (2007), 615--667,
 \href{https://arxiv.org/abs/math.AG/0408337}{arXiv:math.AG/0408337}.

\end{thebibliography}
\end{document}
